% arara: clean: {files: [tfgtfmthesisuam.aux, tfgtfmthesisuam.idx, tfgtfmthesisuam.ilg, tfgtfmthesisuam.ind, tfgtfmthesisuam.bbl, tfgtfmthesisuam.bcf, tfgtfmthesisuam.blg, tfgtfmthesisuam.run.xml, tfgtfmthesisuam.fdb_latexmk, tfgtfmthesisuam.fls, tfgtfmthesisuam.loe, tfgtfmthesisuam.lof, tfgtfmthesisuam.lol, tfgtfmthesisuam.lot, tfgtfmthesisuam.ltb, tfgtfmthesisuam.out, tfgtfmthesisuam.toc, tfgtfmthesisuam.upa, tfgtfmthesisuam.upb, tfgtfmthesisuam.acn, tfgtfmthesisuam.acr, tfgtfmthesisuam.alg, tfgtfmthesisuam.glg, tfgtfmthesisuam.glo, tfgtfmthesisuam.gls, tfgtfmthesisuam.glsdefs, tfgtfmthesisuam.idx,  tfgtfmthesisuam.ilg, tfgtfmthesisuam.xdy, tfgtfmthesisuam.loa, tfgtfmthesisuam.gnuploterrors , tfgtfmthesisuam.mw, tfgtfmthesisuam.fdb_latexmk ]}
% arara: pdflatex: {shell: yes}
% arara: makeglossaries
% arara: makeindex: {style: tfgtfmthesisuam.ist }
% arara: bibtex
% arara: pdflatex: {shell: yes}
% arara: pdflatex: {shell: yes}
% arara: clean: {files: [tfgtfmthesisuam.aux, tfgtfmthesisuam.idx, tfgtfmthesisuam.ilg, tfgtfmthesisuam.ind, tfgtfmthesisuam.bbl, tfgtfmthesisuam.bcf, tfgtfmthesisuam.blg, tfgtfmthesisuam.run.xml, tfgtfmthesisuam.fdb_latexmk, tfgtfmthesisuam.fls, tfgtfmthesisuam.loe, tfgtfmthesisuam.lof, tfgtfmthesisuam.lol, tfgtfmthesisuam.lot, tfgtfmthesisuam.ltb, tfgtfmthesisuam.out, tfgtfmthesisuam.toc, tfgtfmthesisuam.upa, tfgtfmthesisuam.upb, tfgtfmthesisuam.acn, tfgtfmthesisuam.acr, tfgtfmthesisuam.alg, tfgtfmthesisuam.glg, tfgtfmthesisuam.glo, tfgtfmthesisuam.gls, tfgtfmthesisuam.glsdefs, tfgtfmthesisuam.idx,  tfgtfmthesisuam.ilg, tfgtfmthesisuam.xdy, tfgtfmthesisuam.loa, tfgtfmthesisuam.gnuploterrors , tfgtfmthesisuam.mw, tfgtfmthesisuam.fdb_latexmk ]}


\documentclass[epsbased,copyright,final,printable,covers,extendedindex,firstnumbered,tfg,gnuplot]{tfgtfmthesisuam}
\usepackage[spanish]{babel}
\usepackage{xcolor}

\newcommand{\MYhref}[3][blue]{\href{#2}{\color{#1}{#3}}}

\advisor{Pablo Varona}
\levelin{Ingeniería Informática}
\title{Herramientas de tiempo real para la creación de\\ circuitos híbridos entre neuronas vivas y neuronas artificiales}
% \subtitle{Si hace falta subtítulo}
\author{Sergio Hidalgo Gamborino}
\privateaddress{C\textbackslash\ Francisco Tomás y Valiente Nº 11}
\copyrightdate{11 de Febrero de 2025}

\dedication{A mi novia, familia y amigos}
\famouscite{Las neuronas son como las misteriosas mariposas del alma, \\cuyo batir de alas quién sabe si esclarecerá algún día el secreto de la vida mental. \\[0.1em] \begin{flushright}Santiago Ramón y Cajal\end{flushright}}
\prefacefile{inicio/prefacio}
\ackfile{inicio/agradecimientos}
\resumenfile{inicio/resumen}
\abstractfile{inicio/abstract}

\keywords{Biohibrid circuits, Real-time, Neuronal and circuit dynamics, Neural electrophisiology, Operating systems.}
\palabrasclave{Circuitos biohíbridos, Tiempo real, Dinámica neuronal y de circuitos,  Electrofisiología neuronal, Sistemas operativos.}

\coverdata
{
  Escuela Politécnica Superior \\
  Universidad Autónoma de Madrid \\
  C\textbackslash Francisco Tomás y Valiente nº 11
}

\bibliographyconfig{tfgtfmthesisuam}

\datadir{data}
\graphicsdir{img}
\logosdir{img}
\codesdir{codes}


\begin{document}

\chapter{Introducción y Estado del Arte\label{CAP:INTRODUCCIONYESTADOARTE}}{introduccion/introduccion}
  \section{Motivación del trabajo\label{SEC:MOTIVACION}}{introduccion/motivacion}
  \section{Formalización del problema\label{SEC:FORMALIZACION}}{formalizacion/formalizacion}
    \subsection{Modelo neuronal\label{SS:MODELONEURONAL}}{formalizacion/modeloneuronal}
    \subsection{Circuito biohíbrido\label{SS:CIRCUITOBIOHIBRIDO}}{formalizacion/circuitobiohibrido}
    \subsection{Interacción en ciclo cerrado\label{SS:CICLOCERRADO}}{formalizacion/ciclocerrado}
  \section{Estado del Arte\label{SEC:ESTADODELARTE}}{estadoarte/estadoarte}
    \subsection{Sistemas Operativos\label{SS:SSOO}}{estadoarte/ssoo/ssoo}
      \subsubsection{Zephyr\label{SSS:ZEPHYR}}{estadoarte/ssoo/zephyr}
      \subsubsection{RT-Thread\label{SSS:RTTHREAD}}{estadoarte/ssoo/rtthread}
    \subsection{Kernels de tiempo real\label{SS:KERNELS}}{estadoarte/kernels/kernels}
      \subsubsection{RTLinux\label{SSS:RTLINUX}}{estadoarte/kernels/rtlinux}
      \subsubsection{RTAI\label{SSS:RTAI}}{estadoarte/kernels/rtai}
      \subsubsection{Xenomai\label{SSS:XENOMAI}}{estadoarte/kernels/xenomai}
      \subsubsection{Preempt-RT\label{SSS:PTREEMPTRT}}{estadoarte/kernels/preemptrt}
    \subsection{Herramientas para la creación de circuitos biohibridos\label{SS:HERRAMIENTAS}}{estadoarte/herramientas/herramientas}
      \subsubsection{RTHybrid\label{SSS:RTHYBRID}}{estadoarte/herramientas/rthybrid}
      \subsubsection{RTXI\label{SSS:RTXI}}{estadoarte/herramientas/rtxi}
     \subsection{Aplicaciones y nuevas tecnologías\label{SS:APPSYNUEVASTECH}}{estadoarte/appsynuevastech}
   \section{Objetivos\label{SEC:OBJETIVOS}}{introduccion/objetivos}

\chapter{Diseño y desarrollo\label{CAP:DISEÑOEIMPLEMENTACION}}{diseñoeimplementacion/diseñoeimplementacion}
  \section{RTXI para Preempt-RT\label{SEC:RTXIPREEMPTRT}}{diseñoeimplementacion/rtxi}
  \section{Plugins\label{SEC:PLUGINS}}{diseñoeimplementacion/plugins}
  \section{Comprobación del ciclo\label{SEC:COMPROBARCICLO}}{diseñoeimplementacion/comprobarciclo}
  \section{Modelo de Hindmarsh-Rose\label{SEC:HINDMARSHROSE}}{diseñoeimplementacion/hindmarshrose}
  \section{Sinapsis eléctrica\label{SEC:SINAPSISELECTRICA}}{diseñoeimplementacion/sinapsiselectrica}
  \section{Pico Neuron\label{SEC:PICONEURON}}{diseñoeimplementacion/piconeuron}
  \section{Circuito biológico\label{SEC:CIRCBIO}}{diseñoeimplementacion/circuitobiologico}
  
\chapter{Resultados\label{CAP:RESULTADOS}}{resultados}
  \section{Pruebas de rendimiento y latencia\label{SEC:PRUEVASRENDIMIENTOLATENCIA}}{pruebasyrendimiento/pruebasyrendimiento}
    \subsection{RT-Benchmarks\label{SS:RTBENCHMARKS}}{pruebasyrendimiento/rtbenchmarks/rtbenchmarks}
     \subsubsection{Duration\label{SSS:DURATION}}{pruebasyrendimiento/rtbenchmarks/duration}
     \subsubsection{Time Step\label{SSS:TIMESTEP}}{pruebasyrendimiento/rtbenchmarks/timestep}
     \subsubsection{Latency\label{SSS:LATENCY}}{pruebasyrendimiento/rtbenchmarks/latency}
     \subsubsection{Jitter\label{SSS:JITTER}}{pruebasyrendimiento/rtbenchmarks/jitter}
    \subsection{RT-Tests - Cyclictest\label{SS:CYCLICTEST}}{pruebasyrendimiento/cyclictest}
    \subsection{Aislamiento de CPU\label{SS:AISLAMIENTOCPU}}{pruebasyrendimiento/aislamiento}
    \subsection{Resultados y conclusiones\label{SS:RESULTADOS}}{pruebasyrendimiento/resultados}
  \section{Validación con neuronas electrónicas y mediante circuitos biohíbridos\label{SEC:EJEMPLOSCIRCUITOSBIOHIBRIDOS}}{ejemploscircuitos/ejemploscircuitos}
    \subsection{Sinapsis bidireccional entre el modelo de Hindmarsh-Rose y la neurona electrónica\label{SS:HRHRELECTRO}}{ejemploscircuitos/hrhrelectro}
    \subsection{Sinapsis unidireccional entre el modelo de Hindmarsh-Rose y la grabación de neurona real\label{SS:HRREC}}{ejemploscircuitos/hrrec}

\chapter{Conclusiones generales\label{CAP:CONCLUSIONES}}{conclusiones/conclusiones}


\appendix

\chapter{Repositorios\label{CAP:REPOSITORIOSREPO}}{apendices/repositorios/repositorios}
  \section{Útiles y reportes\label{SEC:UTILESYREPORTSREPO}}{apendices/repositorios/utilesyreportes}
  \section{RTXI para Preempt-RT\label{SEC:RTXIPREEMPTREPO}}{apendices/repositorios/rtxipreemptrt}
  \section{Datos y gráficas\label{SEC:DATOSREPO}}{apendices/repositorios/datos}
  \section{RTHybrid para RTXI v3\label{SEC:RTHPRTXI3REPO}}{apendices/repositorios/rthybridrtxplugins}
  \section{Pico Neuron\label{SEC:PICONEURONREPO}}{apendices/repositorios/piconeuron}
  
\end{document}
